\chapter{Caso de estudio}

\section{Ingeniería de requerimientos}

El caso de estudio corresponde al desarrollo incremental de un sistema de punto de venta para farmacia. En la etapa analizada, el problema técnico se concentra en preparar el registro confiable de proveedores y compras recibidas, de modo que una operación comercial futura pueda convertirse en inventario por lote, movimientos trazables y evidencia administrativa.

Los actores funcionales previstos para este flujo son el administrador y el superadministrador, quienes podrán gestionar proveedores y compras, y el vendedor, que no debe operar estos módulos administrativos. La necesidad principal no se limita a registrar datos de una factura o nota de compra: el sistema debe conservar la relación entre proveedor, productos, unidades comerciales, cantidades normalizadas, costos, lote, vencimiento, usuario responsable y estado de la operación. Esta relación es necesaria para que el inventario no nazca de ajustes manuales aislados, sino de una compra recibida y auditada.

El primer corte de implementación del flujo, documentado como Sprint 01 de proveedores y compras recibidas, no entrega todavía la operación completa al usuario final. Su alcance se limita a contratos compartidos, modelo relacional, migración de base de datos y guardrails técnicos. Esta delimitación es relevante para la tesis porque evidencia una estrategia de construcción por capas: antes de implementar endpoints, servicios transaccionales o pantallas, se estabilizan las estructuras de datos y los contratos que sostendrán el flujo.

El segundo corte, documentado como Sprint 02 Backend de Proveedores, transforma esa base en una superficie backend operativa para el directorio de proveedores. Su alcance se limita a listar, consultar, crear y actualizar proveedores mediante endpoints administrativos, con paginación, validación compartida, autorización por rol y auditoría. Este sprint no implementa todavía compras, recepción de mercadería, generación de inventario, anulación de compras ni pantallas de usuario; dichas capacidades permanecen diferidas a cortes posteriores del epic.

El tercer corte, documentado como Sprint 03 Backend de Compras e Inventario, implementa la superficie backend operativa para compras recibidas. El alcance cubre creación y edición de compras en estado \texttt{draft}, listado y consulta de detalle, recepción transaccional, anulación con motivo, generación interna de capas de inventario y movimientos firmados, autorización administrativa y documentación OpenAPI. Este corte no entrega aún frontend de compras, visualización de stock por lote, kardex visual, SIAT, pagos ni cuentas por pagar; esas capacidades permanecen fuera del sprint y deben tratarse como trabajo posterior.

\section{Análisis y diseño}

El diseño mantiene la separación entre contratos compartidos, backend modular y persistencia relacional. En el paquete \texttt{@pharmacy-pos/shared} se definieron esquemas Zod para proveedores, compras, ítems de compra, recepción, anulación y paginación. Estos contratos describen las fronteras de datos esperadas por cliente y servidor, incluyendo paginación \textit{1-based}, tamaño máximo de página, estados de proveedor, estados de compra, fechas puras para \texttt{purchaseDate} y \texttt{expirationDate}, montos con dos decimales y cantidades o factores de conversión con cuatro decimales.

Desde el punto de vista del modelo de datos, la persistencia agrega las entidades \texttt{Supplier}, \texttt{Purchase}, \texttt{PurchaseItem}, \texttt{InventoryBatch} e \texttt{InventoryMovement}. El proveedor conserva razón social, NIT opcional, datos de contacto, estado y relación histórica con compras. La compra registra proveedor, fecha comercial, estado \texttt{draft}, \texttt{received} o \texttt{cancelled}, total, usuario creador, usuario receptor, fechas de recepción o anulación, notas, observaciones de recepción y motivo de anulación. Los ítems guardan producto, unidad comercial, cantidad, costo unitario, lote, vencimiento y capturas de cálculo como \texttt{conversionFactor}, \texttt{baseQuantity}, \texttt{baseUnitCost} y \texttt{lineTotal}.

La decisión de diseño más relevante para inventario es modelar \texttt{InventoryBatch} como una capa interna de inventario originada por un ítem de compra recibido, no como un lote físico único consolidado. Esta estructura permite representar costos distintos para un mismo producto, número de lote y vencimiento, una situación posible cuando varía el costo de reposición. La vista operativa futura podrá agrupar por producto, lote y vencimiento, mientras que la persistencia conserva capas separadas para análisis de costo y trazabilidad.

La trazabilidad se refuerza mediante \texttt{InventoryMovement}. Cada movimiento conserva la capa afectada, el producto, el tipo de movimiento, la cantidad base firmada, el costo unitario base, la referencia de origen, el ítem de referencia, el usuario actor y el motivo. Aunque las reglas que crearán estos movimientos todavía no están implementadas, el diseño de persistencia ya prepara los datos necesarios para recepción y anulación transaccional en sprints posteriores.

Sobre esa base, el Sprint 02 define el módulo backend \texttt{suppliers} como un mini-stack con rutas, controladores, servicios, repositorios y tipos propios. El diseño conserva la separación de responsabilidades: las rutas aplican autenticación y autorización, los controladores validan contratos Zod compartidos y construyen el contexto de auditoría, el servicio concentra reglas de negocio y el repositorio encapsula el acceso a Prisma. Esta estructura evita que la lógica de unicidad, estado histórico o auditoría quede distribuida en la capa HTTP.

Las reglas de proveedor implementadas mantienen el NIT como dato opcional y único cuando existe, permiten conservar proveedores sin NIT y representan la baja operativa mediante el estado \texttt{active} o \texttt{inactive}, sin borrado físico. La lista de proveedores utiliza paginación \textit{1-based}, tamaño de página acotado por contrato compartido, búsqueda por campos administrativos y filtro por estado. La autorización se mantiene alineada con el PRD: \texttt{superadmin} y \texttt{admin} pueden consultar y gestionar proveedores, mientras que \texttt{seller} queda fuera de esta superficie administrativa.

El Sprint 03 extiende el mismo criterio arquitectónico al módulo \texttt{purchases}. El mini-stack de compras contiene rutas, controladores, servicio, repositorio y tipos propios. Las rutas aplican autenticación y autorización por rol; los controladores validan \texttt{PurchasesQuerySchema}, \texttt{CreatePurchaseSchema}, \texttt{UpdatePurchaseSchema}, \texttt{ReceivePurchaseSchema} y \texttt{CancelPurchaseSchema}; el servicio concentra reglas de dominio; y el repositorio encapsula consultas Prisma, transacciones de reemplazo de ítems y escritura de auditoría. Esta división sostiene que la recepción o anulación de una compra no sea una operación HTTP aislada, sino un caso de negocio coordinado.

El inventario relacionado con compras se diseña como un módulo interno, no como una API pública en este sprint. Sus helpers reciben un cliente transaccional desde el servicio de compras y crean o revierten capas \texttt{InventoryBatch} junto con movimientos \texttt{InventoryMovement}. La recepción registra movimientos \texttt{purchase\_received} con cantidad positiva y la anulación de una recepción intacta registra movimientos \texttt{purchase\_cancelled} con cantidad negativa. Esta estrategia permite conservar movimientos firmados y referencias exactas al ítem de compra, al producto, a la capa afectada, al usuario actor y al motivo.

Las reglas de datos quedan alineadas con el objetivo de inventario transaccional. Las compras pueden estar en \texttt{draft}, \texttt{received} o \texttt{cancelled}; las fechas comerciales \texttt{purchaseDate} y \texttt{expirationDate} se manejan como fechas puras; la fecha \texttt{receivedAt} se asigna desde el servidor; los snapshots de conversión, cantidad base, costo base y total de línea se calculan en backend; y los importes se normalizan con escalas diferenciadas para dinero y cantidades. En consecuencia, una compra recibida conserva evidencia histórica aunque luego cambien unidades, costos o datos operativos del producto.

\section{Desarrollo}

El Sprint 01 se desarrolló como base técnica del flujo de proveedores y compras recibidas. La implementación incorporó contratos compartidos en \texttt{packages/shared/src/schemas}, actualizó las exportaciones públicas de \texttt{packages/shared/src/index.ts} y extendió \texttt{backend/prisma/schema.prisma} con los modelos y enumeraciones requeridos. También se generó la migración \texttt{20260511161000\_add\_supplier\_purchase\_inventory\_persistence}, que crea tablas, índices y claves foráneas para los nuevos componentes del dominio.

La migración mantiene la integridad referencial con usuarios, productos y unidades. Por ejemplo, una compra referencia al proveedor y a los usuarios involucrados; un ítem referencia la compra, el producto y la unidad; una capa de inventario referencia el ítem que la origina; y un movimiento referencia la capa, el producto y opcionalmente el usuario actor. Además, se añadieron índices para búsquedas y filtros de proveedores, listados de compras por estado y fecha, consultas por producto, lote, vencimiento, estado de capa y referencias de movimientos.

La unicidad del NIT se trató como una restricción de persistencia compatible con PostgreSQL y Prisma: el campo puede ser nulo para admitir proveedores sin NIT, pero rechaza duplicados cuando el valor existe. Esta decisión responde al contexto operativo donde algunos proveedores pueden no proporcionar información fiscal completa, sin sacrificar control sobre registros duplicados cuando el dato está disponible.

El desarrollo del Sprint 01 no incluye todavía controllers, routes, services, repositories, endpoints HTTP, OpenAPI completo, stores Zustand ni pantallas para proveedores o compras. Tampoco ejecuta aún las reglas de recepción, anulación, creación de capas o creación de movimientos. El resultado de ese sprint debe interpretarse como fundamento arquitectónico y persistente para esas capacidades, no como cierre funcional del epic.

En el Sprint 02 se incorporó el módulo \texttt{backend/src/modules/suppliers} con repositorio, servicio, controlador, rutas y tipos. El repositorio implementa consultas paginadas, búsqueda, filtro por estado, detalle por identificador, creación, actualización y escritura de auditoría. El servicio normaliza campos opcionales, verifica disponibilidad del NIT, transforma valores nulos de persistencia a respuestas de contrato y emite errores de dominio como \texttt{SUPPLIER\_NOT\_FOUND} y \texttt{SUPPLIER\_NIT\_IN\_USE}. Las operaciones de creación y actualización registran eventos \texttt{SUPPLIER\_CREATED} y \texttt{SUPPLIER\_UPDATED} con entidad \texttt{supplier}.

La capa HTTP expone \texttt{GET /api/suppliers}, \texttt{GET /api/suppliers/:id}, \texttt{POST /api/suppliers} y \texttt{PATCH /api/suppliers/:id}. Las rutas aplican \texttt{authenticateRequest} y restringen acceso a \texttt{superadmin} y \texttt{admin}; los controladores usan \texttt{SuppliersQuerySchema}, \texttt{CreateSupplierSchema} y \texttt{UpdateSupplierSchema}. Además, el router principal registra \texttt{/suppliers} y la documentación OpenAPI describe la agrupación de proveedores, parámetros de consulta, esquemas de solicitud y respuesta, seguridad bearer, errores \texttt{401}, \texttt{403}, \texttt{404}, \texttt{409} y validación \texttt{400}. Esta evidencia permite afirmar que el backend de proveedores queda disponible para clientes futuros, sin adelantar todavía una interfaz de usuario.

En el Sprint 03 se incorporó el módulo \texttt{backend/src/modules/purchases} con repositorio, servicio, controlador, rutas y tipos. El repositorio implementa listado paginado con búsqueda, filtros por estado, proveedor y rango de fechas, consulta de detalle con proveedor, usuarios e ítems, creación de borradores, reemplazo transaccional de encabezado e ítems, cambio de estado y auditoría. El servicio valida proveedor activo, usuario autenticado, producto activo, unidad configurada para el producto, ítems no vacíos y duplicados equivalentes. También calcula \texttt{conversionFactor}, \texttt{baseQuantity}, \texttt{baseUnitCost}, \texttt{lineTotal} y \texttt{totalAmount}, evitando depender de cálculos enviados por el cliente.

La recepción de compras se implementó como una transacción única. Antes de recibir, el servicio revalida proveedor, productos, unidades e ítems inventariables; exige lote normalizado y vencimiento no pasado cuando corresponde; y permite ítems no inventariables sin lote ni vencimiento. Dentro de la misma transacción crea capas y movimientos mediante el módulo interno de inventario, marca la compra como \texttt{received}, asigna \texttt{receivedByUserId} y \texttt{receivedAt}, conserva \texttt{receiveNotes} cuando existe y registra auditoría \texttt{PURCHASE\_RECEIVED}. Si falla una regla de dominio, el flujo evita dejar inventario parcial.

La anulación también se implementó con reglas diferenciadas por estado. Una compra en \texttt{draft} puede pasar a \texttt{cancelled} con motivo obligatorio sin afectar inventario. Una compra en \texttt{received} solo puede anularse si las capas originadas por sus ítems permanecen intactas: cantidad disponible igual a cantidad original, estado \texttt{active} y correspondencia con producto, cantidad, costo, lote y vencimiento. Si la validación es satisfactoria, se crean movimientos inversos, se dejan las capas con cantidad disponible cero y estado \texttt{cancelled}, se marca la compra como anulada y se registra \texttt{PURCHASE\_CANCELLED}. Este mecanismo preserva trazabilidad sin borrar historia.

La capa HTTP de compras expone \texttt{GET /api/purchases}, \texttt{GET /api/purchases/:id}, \texttt{POST /api/purchases}, \texttt{PATCH /api/purchases/:id}, \texttt{POST /api/purchases/:id/receive} y \texttt{POST /api/purchases/:id/cancel}. Todas las rutas requieren autenticación y restringen gestión a \texttt{superadmin} y \texttt{admin}; \texttt{seller} queda excluido. La documentación OpenAPI fue actualizada con endpoints, esquemas de compra e ítems, solicitudes de recepción y anulación, parámetros de consulta, seguridad bearer y errores esperados de validación, autenticación, autorización, no encontrado y conflictos de dominio.

\section{Pruebas}

La verificación técnica registrada para el Sprint 01 se enfocó en artefactos de contratos y persistencia. Se ejecutó correctamente la generación del cliente Prisma mediante \texttt{pnpm --filter @pharmacy-pos/backend prisma:generate}. También se ejecutó correctamente el typecheck del paquete compartido con \texttt{pnpm --filter @pharmacy-pos/shared typecheck}, lo que respalda que los contratos exportados compilan como parte del monorepo.

El typecheck del backend fue ejecutado, pero permanece bloqueado por una deuda previa no relacionada con el sprint: una incompatibilidad de tipos en el módulo de autenticación respecto al estado de usuario \texttt{blocked}. Esta falla no proviene de los modelos de proveedores, compras ni persistencia de inventario introducidos en el Sprint 01. La evidencia documentada distingue, por tanto, entre validación satisfactoria del alcance compartido y una deuda técnica existente fuera del corte implementado.

La revisión de la migración confirmó que el cambio crea estructuras nuevas y no incluye operaciones destructivas sobre tablas existentes, como borrado de tablas, truncamiento de datos o eliminación de columnas. No se realizó QA manual en navegador porque el sprint no incorporó rutas, pantallas ni flujos web.

Para el Sprint 02, la evidencia de cierre registrada indica que el QA manual de API fue solicitado, pero no pudo completarse porque el host local rechazó la conexión al intentar autenticar el usuario seed en \texttt{POST http://localhost:4000/api/auth/login}. Como consecuencia, no se ejercitaron los endpoints de proveedores ni los casos negativos de autorización o conflicto desde una sesión real. La documentación del sprint consigna este bloqueo y mantiene separada la existencia del backend implementado respecto de la verificación manual pendiente por disponibilidad del servidor local.

Para el Sprint 03, la estrategia de validación se apoya principalmente en contratos compartidos, validaciones de dominio en servicios, transacciones Prisma, autorización por roles, auditoría y documentación explícita de errores. La evidencia implementada muestra que la recepción y la anulación se ejecutan desde el servicio de compras con un cliente transaccional compartido por los helpers de inventario, lo que reduce el riesgo de estados parciales entre compra, capas, movimientos y auditoría. También se registran códigos de error específicos para estados inválidos, proveedor o producto no válido, unidad no configurada, ítems duplicados, vencimiento inválido, motivo faltante y capas no reversibles.

El QA manual del Sprint 03 fue solicitado para API/backend, pero quedó bloqueado por la misma condición ambiental registrada en el ticket de QA: no existía un listener activo en los puertos locales esperados y la autenticación contra \texttt{POST http://localhost:4000/api/auth/login} fue rechazada por el host local. Por ello, no se afirma que los endpoints hayan sido ejercitados manualmente en ejecución local durante ese cierre. La tesis toma como evidencia de este sprint el código implementado, la documentación OpenAPI y la trazabilidad del sprint, dejando la verificación manual como una limitación operativa pendiente de disponibilidad del servidor.

\section{Estimación de costos}

En esta sección se presenta una estimación de costos del caso de estudio, considerando recursos humanos, tecnológicos, operativos y otros elementos pertinentes.
